% Proposed B.153; owned source, no central edits.
\subsection*{The actual compact stabilizer and a geometric curve}
\statement{B.153}{Proved}{The two actual completed sections intersect in the procyclic closure of the marked inner source}
\label{stmt:B153}
Use B.148's actual coarse group $E^0=T_0\rtimes S$, with
$T_0=F(r,\bar r,w)$, $S=F(m,u)$, and graph section
$H=\{h_s^{-1}s:s\in S\}$. Both section completions embed in the
actual full completion of $E^0$, and
\begin{equation}\label{eq:B153-conclusion}
 \widehat H\cap\widehat S
 =\{s\in\widehat S:\alpha_s(r)=r\}
 =\overline{\langle m\rangle}
 =\widehat{\mathbb Z}(m).
\end{equation}
The intersection is taken in $\widehat{E^0}$; no completion of an
abstract subgroup intersection is presumed.
\begin{proof}
\emph{The actual completed return equation.}
Every finite quotient of the finite-rank free fiber $T_0$ is refined by
a characteristic finite-index subgroup. Its finite automorphism action
extends it to a finite quotient of the split group. Thus $T_0$ has full
induced topology, and
$\widehat{E^0}=\widehat T_0\rtimes\widehat S$.
The source $S$ and the graph section $H$ each project isomorphically to
$S$, so both have full topology as well. Their actual sections and the
ordered cocycle extend continuously. The identities
\[
 \alpha_s(r)=h_srh_s^{-1},\qquad\chi_x(h_s)=0,
 \qquad\chi_x(r)=1
\]
therefore hold for every compact source. Since $r$ is a free basis
letter, the procyclic free factor in
$\widehat T_0=\widehat{\mathbb Z}(r)\amalg\widehat F(\bar r,w)$
has $C_{\widehat T_0}(r)=\widehat{\mathbb Z}(r)$.
If $s$ fixes $r$, write $h_s=r^\lambda$ and apply $\chi_x$ to obtain
$\lambda=0$. This proves the first equality in
\eqref{eq:B153-conclusion}. All sources commute with the actual
involution $\tau$, so such a source fixes $\bar r=\tau r\tau$ too.

\emph{The complete geometric marking.}
Put $f=w^{-1}\bar r^{-1}$ and $g=r^{-1}w$. The identity
$rgf\bar r=1$ identifies the actual $T_0$ with the fundamental group
of an oriented four-punctured sphere with positive peripheral tuple
$(r,g,f,\bar r)$. In this precise tuple, B.148 gives
\begin{equation}\label{eq:B153-geometric-actions}
\begin{gathered}
 m(r)=r,\quad m(\bar r)=\bar r,\quad
 m(g)=f,\quad m(f)=f^{-1}gf,\\
 u(r)=w^{-1}rw,\quad u(g)=w^{-1}gw,\quad
 u(f)=wfw^{-1},\quad u(\bar r)=w\bar r w^{-1},
 \qquad w=rg.
\end{gathered}
\end{equation}
Thus $m$ is the inverse Hurwitz half-twist on $(g,f)$. Modulo the
global inner automorphism by $w^{-1}$, $u$ conjugates the first pair
$(r,g)$ by $w^{-2}$ and fixes the other pair; its outer class is a
squared separating twist. These are actual orientation-preserving
mapping classes. The simple curve $\gamma$ separating the two pairs
is represented up to conjugacy and inversion by $r\bar r$, because
$\bar r r=(gf)^{-1}$.

\emph{The finite-index source and the pure subgroup.}
The pillowcase comparison gives the homomorphism from the full
four-punctured-sphere mapping-class group $M$ to
$\mathrm{PSL}_2(\mathbb Z)$ with Klein-four kernel, and identifies
curve slopes via the branched torus cover; see the complete proofs in
\cite[\S2.2.5, pp.~57--60, Propositions~2.6--2.7]{CEFarbMargalit2011}.
Our homology basis is exactly the translation basis $(z,w)$ of B.149,
where $z=\bar r r$. The source matrices are
\[
 M_0=\begin{pmatrix}1&-1\\0&1\end{pmatrix},\qquad
 U_0=\begin{pmatrix}1&0\\4&1\end{pmatrix}.
\]
B.149 proves their projective faithfulness on all of $S$, so the actual
outer source has neither an inner nor a Klein-four kernel.
Conjugating by $\operatorname{diag}(2,1)$ gives the upper $-2$ and
lower $2$ generators of $\Gamma(2)$ modulo signs. To verify generation,
reduce a primitive column $(a,c)$ with $a$ odd and $c$ even by adding
even multiples of either coordinate to the other. The larger absolute
value strictly decreases: equality with the smaller one is excluded
by parity. This ends at $(\pm1,0)$, after which a sign and an even
upper shear suffice. Entrywise conjugation therefore proves
\[
 \rho(S)=\Gamma_0(4)/\{\pm I\},\qquad
 \operatorname{diag}(2,1)\Gamma_0(4)\operatorname{diag}(2,1)^{-1}
 =\Gamma(2).
\]
The first subgroup has index six in $\mathrm{PSL}_2(\mathbb Z)$:
it stabilizes $(1,0)$ in the six-element projective line over
$\mathbb Z/4$, on which the standard shear and quarter-turn act
transitively. Hence $S$ has index twenty-four in $M$.

For the primary profinite theorem use only the pure group $M_p$.
The actual puncture permutation gives
$S_p=\ker(\operatorname{exp}_m\bmod2)$.
Its image is $(\Gamma_0(4)\cap\Gamma(2))/\{\pm I\}$, of index two in
$\Gamma(2)/\{\pm I\}\cong M_p$; the Klein-four kernel meets $M_p$
trivially. If a source preserves $\gamma$, its matrix sends the slope
$\pm(1,0)^t$ to itself. It is therefore, up to sign, a power of $M_0$.
Projective faithfulness and the literal $m(r\bar r)=r\bar r$ prove
\begin{equation}\label{eq:B153-discrete-stabilizer}
 S\cap\operatorname{Stab}_M(\gamma)=\langle m\rangle,
 \qquad S_p\cap\operatorname{Stab}_{M_p}(\gamma)=\langle m^2\rangle.
\end{equation}

\emph{The full completion and geometric stabilizer.}
The genus-zero congruence subgroup property identifies
$\widehat M_p$ with its closure in
$\operatorname{Out}(\widehat{\pi_1(S_{0,4})})$; this is
\cite{Asada2001}, with an independent full proof and the explicit
$n=4$ deduction in \cite{BMcReynoldsPureBraidCSP2012}.
Boggi's Proposition~3.4 and Theorem~4.2 identify the procongruence
curve completion with its realization by conjugacy classes of
profinite simple curves and give the stabilizer of a geometric vertex
as the closure of its discrete stabilizer
\cite{CEBoggiProcongruence2014}.
The hypotheses include our separating essential curve on $S_{0,4}$.
Applying this statement to the open subgroup $\widehat S_p$ gives
\[
 \operatorname{Stab}_{\widehat S_p}(\gamma)
       =\overline{\langle m^2\rangle}.
\]
Here the primary theorem uses conjugacy classes up to inversion, so
fixing the actual word $r\bar r$ suffices.

For precision in its cofinal-cover input, one may begin inside the
characteristic subgroup $K_3=[\pi_1,\pi_1]\pi_1^3$. Its index is
$27$, all four peripheral orders are $3$, and the cover has $36$
punctures and Euler characteristic $-54$, hence compact genus $10$.
Further covers compactify to possibly branched covers of this surface
and have genus at least two. Intersecting any prescribed characteristic
subgroup with $K_3$ preserves cofinality before applying the
homology-cover refinement of
\cite[Lemma~3.10, Theorem~7.7 and Corollary~7.8]{BBoggiGaloisCoverings2014}.
This avoids using a small-cover genus assertion outside its range.
The complete original separating-curve and cover dependencies are
recorded in the independent primary audit below.

A compact source fixing $r$ also fixes $\bar r$, hence $\gamma$.
If it lies in $\widehat S_p$, it belongs to
$\overline{\langle m^2\rangle}$. Otherwise multiply it by $m^{-1}$
and apply the even case. Conversely all profinite powers of $m$ fix
$r$. This proves \eqref{eq:B153-conclusion} in the actual completion.
\end{proof}
The complete primary/application audit is
\path{work_DE/next_marginal_source/Boggi_S04_primary_marking_audit31.md};
the exact geometric proof is
\path{work_CE/compact_one_family_stabilizer_geometric31.txt}.
This uses full genus-zero congruence comparison, not faithfulness of an
arithmetic matrix completion. It does not decide remaining contact
fixed-kernel generation or the original full-base B.75 criterion.
